\chapter*{Introduction générale}
\addcontentsline{toc}{chapter}{Introduction générale}

Les entreprises reposent aujourd'hui de plus en plus sur des infrastructures
numériques pour assurer la continuité de leurs services. Au delà de la
simple mise en ligne d'une application, il s'agit de pouvoir la déployer,
la surveiller, en contrôler l'accès et réagir rapidement à un incident,
sans dépendre d'interventions manuelles dispersées sur chaque serveur.
C'est dans cette optique qu'un stage a été effectué au sein d'AG
Technologies, entreprise camerounaise spécialisée dans la conception de
solutions digitales (présentée en détail au chapitre~\ref{chap:entreprise}),
du 06 juillet 2026 au 05 août 2026. La mission confiée portait initialement
sur la prise en main des outils internes de l'entreprise et l'amélioration
de projets logiciels déjà existants, distincts du projet finalement au
coeur de ce rapport.

\section*{Contexte et problématique}

AG Technologies exploite un parc de serveurs physiques hétérogènes
hébergeant des applications critiques pour son activité, parmi lesquelles
AGT-BOT, AGT ERP et MboaPay. Avant le développement du projet présenté
dans ce rapport, ce parc était piloté de façon manuelle et dispersée
(section~\ref{chap:entreprise} du chapitre suivant) : brancher un serveur,
y lancer des conteneurs, surveiller sa charge, déployer une application,
l'exposer publiquement, gérer les accès des personnes intervenant dessus,
et ajuster ses ressources en cas de montée en charge constituaient autant
d'opérations menées séparément, sans point de contrôle unique ni
traçabilité homogène des actions effectuées.

Cette dispersion posait un problème de fond : plus le nombre de serveurs
et d'applications augmentait, plus le risque d'erreur, d'oubli ou d'accès
non maîtrisé augmentait avec lui, sans qu'aucun outil interne ne permette
d'avoir une vue d'ensemble ni un contrôle centralisé. C'est ce constat qui
a conduit à la définition d'un projet interne, nommé AGT Infra, présenté
en détail dans la suite de ce rapport.

\section*{Objectifs du projet}

L'objectif général du projet est de construire une interface interne
unique capable de piloter l'essentiel du parc de serveurs depuis un seul
point d'entrée, sans réimplémenter les outils déjà éprouvés sur lesquels
elle s'appuie (Docker, SSH, Nginx, Certbot, GitHub) : AGT Infra se
positionne comme un chef d'orchestre de ces briques, pas comme un
remplacement. Cet objectif général se décline en plusieurs objectifs
spécifiques :

\begin{itemize}
  \item Centraliser la gestion des serveurs physiques et des conteneurs
        Docker déjà en place, avec un inventaire unique et une connectivité
        exclusivement fondée sur SSH.
  \item Automatiser le déploiement continu des applications à partir de
        dépôts Git, déclenché par des évènements applicatifs (webhooks)
        vérifiés.
  \item Permettre l'exposition sécurisée, sous nom de domaine et en HTTPS,
        des applications déjà déployées.
  \item Superviser l'état des serveurs, à la demande comme en continu, et
        déclencher des alertes en cas de dépassement de seuils.
  \item Mettre en place un contrôle d'accès dynamique, fondé sur des rôles
        et des permissions pilotés en base de données plutôt que figés
        dans le code.
  \item Tracer systématiquement toute action sensible dans un journal
        d'audit horodaté et attribué.
  \item Fournir une console web intégrée à l'interface, permettant
        d'intervenir directement sur un serveur ou un conteneur sans ouvrir
        de session SSH séparée.
\end{itemize}

Certains besoins identifiés dès la définition du projet, comme les
sauvegardes automatiques des bases de données ou la montée en charge
automatique, restent hors du périmètre effectivement réalisé à ce jour ;
leur état d'avancement est précisé au fil des chapitres concernés. La
gestion des machines virtuelles et de l'hyperviseur associé n'entre quant
à elle pas dans le périmètre du projet, qui se concentre exclusivement sur
le pilotage de serveurs physiques et de conteneurs Docker déjà en place.

\section*{Structure du rapport}

Le présent rapport est organisé de la façon suivante :

\begin{itemize}
  \item \textbf{Chapitre~\ref{chap:entreprise} : Présentation de
        l'entreprise AGT} : activité et organisation d'AG Technologies,
        puis environnement technique dont elle disposait avant le
        développement d'AGT Infra.
  \item \textbf{Chapitre~\ref{chap:etat-art} : État de l'art} : panorama
        des solutions existantes du marché et justification du choix d'un
        développement interne.
  \item \textbf{Chapitre~\ref{chap:etude-technique} : Analyse et
        conception} : analyse des besoins (acteurs, exigences,
        diagramme de cas d'utilisation), architecture générale de
        l'outil, modélisation de ses données et comportements clés.
  \item \textbf{Chapitre~\ref{chap:realisation} : Implémentation} :
        environnement et outils de développement, méthodologie de
        travail suivie, réalisation du projet module par module,
        difficultés rencontrées et déploiement.
  \item \textbf{Chapitre~\ref{chap:resultats} : Résultats et bilan} :
        résultats obtenus, workflows et captures d'écran de
        l'application, tests effectués, limites actuelles et
        perspectives d'évolution du projet.
  \item \textbf{Conclusion générale} : bilan des compétences techniques
        et humaines acquises durant le stage, et apport de cette
        expérience.
\end{itemize}